\chapter{「アライメント」タブの設定}\label{chap:align}

アライメントタブの設定は、\textbf{お手本（参照画像）の作り方}と、\textbf{位置合わせ領域の置き方・合わせ方}を決めます。
ここを変えたら \ui{再アライメント} を押します（品質評価はやり直しません）。

% ------------------------------------------------------------------------------
\section{アライメント（位置合わせ）}

\sidebyside[0.5]{insp_align}{%
  \mk[e]{method}{1}\mk[e]{mode}{2}\mk[e]{apsize}{3}\mk[e]{search}{4}\mk[e]{placement}{5}\mk[e]{ref}{6}\mk[e]{refine}{7}\mk[e]{report}{8}}{%
  \begin{marklist}
    \item 位置合わせ方法
    \item 対象モード
    \item 位置合わせ領域の大きさ
    \item 位置ずれの探索範囲
    \item 位置合わせ領域の配置
    \item 参照フレーム（品質上位 \%）
    \item 参照の反復精密化（2パス）
    \item \textbf{結果の内訳}：アライメントが終わると出ます。
  \end{marklist}
  \smallskip
  \badge{8} には、自動で判定した対象モードとお手本にしたフレームの番号、除外したフレームの内訳、
  位置合わせ領域の数と大きさ、うまく合わなかった所をどう扱ったかが出ます。結果が変なときの手がかりになります。}

\begin{setting}{\badge{1} 位置合わせ方法}{複数領域の局所位置合わせ（推奨）}{複数領域の局所位置合わせ／画像全体の位置合わせのみ（高速）}
  \begin{itemize}[itemsep=0pt, topsep=2pt]
    \item \uil{複数領域の局所位置合わせ（推奨）}：画像を位置合わせ領域に分け、領域ごとにずれを合わせ、領域ごとに写りのよいフレームを選びます。
          大気のゆらぎによる場所ごとのゆがみまで直せるので、いちばんシャープになります。
    \item \uil{画像全体の位置合わせのみ（高速）}：画像全体を1つとして平行移動だけで合わせます。速いですが、場所ごとのゆがみは残ります。
          このときは、参照に使う割合（オレンジの線）がそのままスタックの採用率になります。
  \end{itemize}
  \meyasu{ふだんは \uil{複数領域}。とにかく早く結果を見たいとき、ゆらぎがとても小さい日には \uil{画像全体のみ} も使えます。}
\end{setting}

\begin{setting}{\badge{2} 対象モード}{自動}{自動／惑星／月・太陽}
  全体の位置合わせのやり方を、対象に合わせて選びます。
  \begin{itemize}[itemsep=0pt, topsep=2pt]
    \item \uil{惑星}：暗い背景の中の明るい円盤として扱い、その重心で大まかに合わせてから精密に合わせます。惑星が画面の中を動いても追いかけられます。
    \item \uil{月・太陽}：画面いっぱいの模様として、画像全体の模様の一致で合わせます。
  \end{itemize}
  \meyasu{\uil{自動} のままで、画像から判定されます（判定結果は \badge{8} に出ます）。月の欠け際のように、暗い部分が大きい月面で
  \uil{惑星} と判定されてうまくいかないときは、\uil{月・太陽} を選び直してください。}
\end{setting}

\begin{setting}{\badge{3} 位置合わせ領域の大きさ}{自動}{自動／32／48／64／96／128／200 px}
  位置合わせ領域（四角）1つの一辺の画素数です。\uil{自動} では、対象の明るい部分の広がりから決めます（例：448 画素の木星では 48 画素、
  1024×768 の月面では 128 画素）。隣どうしの四角は半分ずつ重なるように並びます。
  \begin{itemize}[itemsep=0pt, topsep=2pt]
    \item \textbf{小さくする}と、細かい場所ごとのゆがみに追いつけますが、四角の中の模様が少なくなり、ノイズで合わせ間違えやすくなります。
    \item \textbf{大きくする}と、合わせ方は安定しますが、場所ごとのゆがみの補正は粗くなります。
  \end{itemize}
  \meyasu{まず \uil{自動}。ゆらぎが大きい日で、スタック結果の一部がぼやけるときは1段小さく、模様の少ない対象（金星・火星のうす模様）で
  結果がゆがむときは1段大きくします。惑星の円盤に数十〜百数十個並ぶくらいが目安です。}
\end{setting}

\begin{setting}{\badge{4} 位置ずれの探索範囲}{±16 px}{±8〜±32 px}
  位置合わせ領域ごとに、お手本からどこまでずれた所を探すかです。
  \meyasu{ゆらぎが大きく、像が大きく踊る日は ±24〜±32。小さい像やゆらぎの小さい日は ±8〜±12 にすると速く、合わせ間違いも減ります。
  結果の内訳の「一致度不足で補間」が多いときは、範囲を広げると改善することがあります。}
\end{setting}

\begin{setting}{\badge{5} 位置合わせ領域の配置}{自動配置}{自動配置に戻す／すべて消去}
  \ui{自動配置に戻す} で、手で足したり消したりした配置を自動の配置に戻します。\ui{すべて消去} で全部消し、自分で置き直せます（第\ref{sec:apedit}節）。
\end{setting}

\begin{setting}{\badge{6} 参照フレーム（品質上位 \%）}{25 \%}{5〜50 \%}
  品質の上位何\%のフレームを重ねてお手本を作るかです。品質グラフのオレンジの線と連動します。
  \meyasu{お手本はくっきりしていてノイズが少ないほど合わせやすくなります。フレームが数千枚あれば 10〜25\%、
  数百枚しかなければ 25〜50\% にして、お手本のノイズを減らします。}
\end{setting}

\begin{setting}{\badge{7} 参照の反復精密化（2パス）}{オン}{オン／オフ}
  一度位置を合わせて作った結果を新しいお手本にして、もう一度合わせ直します。お手本がよりシャープになり、位置合わせの精度が上がります。
  \meyasu{オンのまま。アライメントの時間を短くしたいときだけオフにします。}
\end{setting}

\begin{tip}[結果の内訳の読み方]
  \begin{itemize}[itemsep=0pt, topsep=2pt]
    \item \uil{除外：類似度不足 / 位置ずれ超過 / 外れ値}：品質評価タブの詳細設定（第\ref{chap:quality}章）の基準で除外した枚数です。
    \item \uil{一致度不足で補間}：模様がはっきりせず、うまく合わなかった位置合わせ領域の割合。その領域は周りの領域のずれから補って使います。
          惑星の縁や模様の少ない所では、ある程度出るのが普通です。
    \item \uil{近傍との差で制限}：隣の領域と比べて不自然に大きくずれた結果を抑えた割合。
    \item \uil{共通の位置ずれへ退避}：1枚のフレームの中で領域ごとのずれがばらばらで信用できないとき、そのフレームは全体に共通のずれだけで合わせます。
          木星の縞のように模様が一方向に並んだ対象では、この割合が高くなることがあります。
  \end{itemize}
\end{tip}

% ------------------------------------------------------------------------------
\section{詳細設定（位置合わせ領域）}

\sidebyside[0.5]{insp_aadv}{\mk[w]{minscore}{1}\mk[w]{gradient}{2}\mk[w]{level}{3}}{%
  \begin{marklist}
    \item 一致度の下限（下回ると補間）
    \item 配置する模様の強さ（比）
    \item 配置する明るさ（比）
  \end{marklist}}

\begin{setting}{\badge{1} 一致度の下限（下回ると補間）}{0.5}{0〜1}
  位置合わせ領域ごとの「お手本との一致度」（1 が完全に一致）がこの値より低いときは、その結果を使わず、周りの領域のずれから補います。
  \meyasu{合わせ間違いでスタック結果の一部が崩れるときは 0.6〜0.7 に上げます。補間が多すぎる（内訳の「一致度不足で補間」が大きい）ときは 0.3〜0.4 に下げます。}
\end{setting}

\begin{setting}{\badge{2}\,\badge{3} 配置する模様の強さ・明るさ（比）}{0.6・0.15}{0〜1 くらい}
  自動配置で、位置合わせ領域を置く場所の条件です。画像の中でいちばん強い模様・いちばん明るい所に対する割合で決めます。
  模様が弱すぎる所や暗すぎる所（背景の宇宙など）には置きません。
  \meyasu{月の欠け際の暗い所や、模様の乏しい所にも置きたいときは、模様の強さを 0.3〜0.4、明るさを 0.05〜0.1 に下げます。
  背景のノイズの上にまで置かれてしまうときは上げます。}
\end{setting}

% ------------------------------------------------------------------------------
\section{位置合わせ領域を手で置く}\label{sec:apedit}

自動の配置が気に入らないときは、手で足したり消したりできます。

\begin{enumerate}[itemsep=2pt]
  \item プレビューの下の \uil{配置を編集} をオンにします。
  \item プレビューを\textbf{クリック}すると、その場所に位置合わせ領域が1つ増えます。
  \item 領域をクリックして選び、\key{delete} を押すと消えます。
  \item 間違えたら \key{⌘}\key{Z} で取り消し、\key{shift}\key{⌘}\key{Z} でやり直しです。
  \item 配置を変えたら \ui{再アライメント} を押します。
\end{enumerate}

\begin{caution}
  手で置いた配置は、\ui{自動配置に戻す} を押すまで使われます。別のファイルを開いたときは自動の配置に戻ります。
\end{caution}

% ------------------------------------------------------------------------------
\section{品質で色分けする}\label{sec:heatmap}

プレビューの下の \ui{表示} メニューで \uil{品質で色分け} を選ぶと、位置合わせ領域の枠が、その場所の品質の平均で色分けされます。
青いほど品質が低く、緑・黄・赤になるほど高い場所です。「どこが最後までゆらぎに負けていたか」が分かります。

\screen[0.72\linewidth]{heatmap}{品質で色分けした位置合わせ領域}{\mk[ne]{legend}{1}}

\begin{marklist}
  \item 色の見本（左が品質が低い、右が高い）。
\end{marklist}
